% Figure IV.14 --- cartel stability is a two-lever phase diagram.
%
% Reader question: under which detection probabilities and sampled losses
% does collusion stop paying?
% Claim: the repeated-game boundary $V_C=\pi_D$ separates a region where
% deviation dominates from one where the discounted collusive stream still
% wins, and a rare severe loss can sit on the same boundary as a frequent
% smaller one.
% Evidence: $V_C=(\pi_C-p_dL)/(1-\delta(1-p_d))$, the two worked points
% (rare/severe; frequent/smaller) that this section marks on the boundary.
% Must distinguish: the two named regions (a verdict each) from the two
% specific worked points on the boundary between them.
% Grammar chosen: payoff inequality as a phase/regime plot (atlas first
% choice for IV/fig:cartel-game-inline). Grammar rejected: a decorative
% seesaw or five ticks with no scale (atlas reject list) -- neither carries
% the boundary's shape or the two worked points.
% Five-second test: a reader can find both worked points on the same
% boundary curve and say which region each one falls in.
\begin{figure}[H]
\centering
\pdfigurehue{pdgold}%
\begin{tikzpicture}[pd figure,x=1cm,y=1cm]
  \draw[pd rule,->] (-0.30,-1.45) -- (5.30,-1.45);
  \draw[pd rule,->] (-0.30,-1.45) -- (-0.30,3.15);
  \node[pd label,anchor=west] at (5.45,-1.45) {detection probability $p_d$};
  \node[pd label,rotate=90,anchor=south] at (-0.70,.80) {loss if detected $L$};
  % ---- the two regions, edged by the boundary they share -------------------
  \begin{scope}[on background layer]
    \fill[pd ready fill]
      (0.30,2.65) -- (1.70,2.00) -- (3.10,1.05) -- (3.95,.35) -- (4.50,-.15)
      -- (4.90,-.85) -- (4.90,3.15) -- (0.30,3.15) -- cycle;
    \fill[pd warn fill]
      (0.30,2.65) -- (1.70,2.00) -- (3.10,1.05) -- (3.95,.35) -- (4.50,-.15)
      -- (4.90,-.85) -- (4.90,-1.45) -- (0.30,-1.45) -- cycle;
  \end{scope}
  \draw[pd focus rule,line width=1.4pt]
    (0.30,2.65) -- (1.70,2.00) -- (3.10,1.05) -- (3.95,.35) -- (4.50,-.15)
    -- (4.90,-.85);
  \node[pd focus label,anchor=south west] at (0.45,2.72)
    {boundary $V_C=\pi_D$};
  \node[pd label,align=center,text width=2.5cm] at (3.60,2.55)
    {deviation pays\\cartel unstable};
  \node[pd label,align=center,text width=2.3cm] at (1.72,-.05)
    {collusion pays\\discounted stream wins};
  % ---- the two worked points, on the shared boundary ------------------------
  \node[pd datum] at (1.70,2.00) {};
  \node[pd tag,anchor=south east] at (1.53,2.10) {rare, severe};
  \node[pd datum] at (3.95,.35) {};
  \node[pd tag,anchor=north east] at (3.78,.20) {frequent, smaller};
\end{tikzpicture}
\caption{\textbf{Cartel sustainability in the detection--loss plane.} The
curve marks $V_C=\pi_D$: below it collusion pays; above it deviation pays.
The two marked policies share the boundary, so detection frequency alone does
not determine stability. [internal]}
\label{fig:cartel-game-inline}
\end{figure}
